\section{Execution: what readiness and placement cost}
\label{sec:execution}

\Cref{sec:protocol} states what must happen before what. This section states
what that costs in wall time, where the tasks may run, and which history each
window ranges over when the packages live in several repositories. The duration
propositions are informative in the sense of \spec{13.11}: they bound
durations of schedules that respect \cref{def:readiness}, and no conforming
implementation is required to attain them. The two optional profiles of
\spec{28} and \spec{27} are transcribed as definitions, with the one invariance
each promises, in \cref{sec:exec-distributed,sec:exec-polyrepo}. Nothing here
extends the theorems of \cref{sec:guarantees}, which are stated for a single
history.

\subsection{Build readiness durations}
\label{sec:exec-durations}

\spec{19.2a} promises a happens-before relation and no duration; \spec{13.11}
nevertheless says what each relation buys when nothing else limits the run.
Throughout, package $i$ has a build of duration $b_i$ and a publication of
duration $p_i$, both fixed; a task starts as soon as its predecessors in the
task graph of \cref{def:taskgraph} have finished and a slot is free; and one
relation is declared on every edge.

\begin{proposition}[the longest path, \spec{13.11}]
\label{prop:exec-longest}
With unlimited build and publication slots a run lasts exactly the longest
path $L$ of its task graph, measured in task durations. A weaker relation
never lengthens it.
\end{proposition}

\begin{proof}
The task graph is acyclic (\cref{prop:readiness-acyclic}), so its tasks can be
taken in a topological order. With unlimited slots a task starts at the
greatest finish time of its predecessors, or at $0$ if it has none, and by
induction along that order its finish time is the longest path ending at it.
The run ends at the greatest finish time, which is the longest path of the
graph. A weaker relation contributes a subset of the edges
(\cref{prop:readiness-acyclic}), so every path of its task graph is a path of
the stronger one's, and the longest is no longer.
\end{proof}

\begin{proposition}[a chain, \spec{13.11}]
\label{prop:exec-chain}
Let packages $1, \dots, n$ form a chain in which package $k+1$ consumes package
$k$, under one relation throughout. With unlimited slots the run lasts
\[
  \sum_{i=1}^{n} (b_i + p_i) \quad\text{under \code{publish}},\qquad
  \max_{1 \le k \le n}\Bigl(\sum_{i=1}^{k} b_i + \sum_{i=k}^{n} p_i\Bigr)
  \quad\text{under \code{build}},
\]
\[
  \max_{1 \le k \le n}\Bigl(b_k + \sum_{i=k}^{n} p_i\Bigr)
  \quad\text{under \code{none}}.
\]
\end{proposition}

\begin{proof}
Under \code{publish}, $\exbuild{k+1}$ waits for $\expub{k}$, so the $2n$ tasks
form one path and the sum is its length. Under \code{build}, the builds form
the path $\exbuild{1} \to \dots \to \exbuild{n}$ and $\exbuild{k}$ ends at
$B_k = \sum_{i \le k} b_i$; the publication of $k$ waits for $\exbuild{k}$ and
for $\expub{k-1}$, so its finish $E_k$ satisfies $E_k = \max(B_k, E_{k-1}) +
p_k$ with $E_0 = 0$. By induction $E_k = \max_{j \le k}\bigl(B_j +
\sum_{i=j}^{k} p_i\bigr)$: the case $k = 1$ is $B_1 + p_1$, and the step is
$\max\bigl(B_k + p_k,\ \max_{j < k}(B_j + \sum_{i=j}^{k-1} p_i) + p_k\bigr)$.
Under \code{none} every build starts at $0$ and $\exbuild{k}$ ends at $b_k$;
the same recurrence with $b_k$ in place of $B_k$ gives the third form. In each
case the run lasts $E_n$ by \cref{prop:exec-longest}.
\end{proof}

\begin{corollary}[equal durations]
\label{cor:exec-equal}
With $b_i = b$ and $p_i = p$ for all $i$, the three forms are $n(b+p)$,
$\max(nb + p,\ b + np)$ and $b + np$. \code{build} saves $(n-1)\min(b,p)$
over \code{publish}, and \code{none} saves a further $(n-1)\max(b-p, 0)$, so
the second saving exists only where builds outlast publications.
\end{corollary}

\begin{proof}
Under \code{build} the maximand $kb + (n-k+1)p$ is affine in $k$, so the
maximum is at $k = 1$ or $k = n$. Under \code{none} the maximand $b + (n-k+1)p$
is greatest at $k = 1$. The differences follow by cases on $b \ge p$.
\end{proof}

\begin{proposition}[fan-out and fan-in, \spec{13.11}]
\label{prop:exec-fan}
For one provider $0$ and consumers $i \ge 1$, each consuming $0$ and nothing
else, the run lasts
\[
  b_0 + p_0 + \max_i (b_i + p_i),\qquad
  b_0 + \max_i\bigl(\max(b_i, p_0) + p_i\bigr),\qquad
  \max_i\bigl(\max(b_i,\ b_0 + p_0) + p_i\bigr)
\]
under \code{publish}, \code{build} and \code{none} respectively. For providers
$i \ge 1$ of one consumer $c$, each consumed by $c$ and nothing else, it lasts
\[
  \max_i (b_i + p_i) + b_c + p_c,\qquad
  \max\bigl(\max_i b_i + b_c,\ \max_i (b_i + p_i)\bigr) + p_c,\qquad
  \max\bigl(b_c,\ \max_i (b_i + p_i)\bigr) + p_c.
\]
\end{proposition}

\begin{proof}
Fan-out. Under \code{publish}, consumer $i$'s build starts at $b_0 + p_0$ and
its publication follows its build, giving $b_0 + p_0 + b_i + p_i$; take the
maximum. Under \code{build}, its build starts at $b_0$ and ends at $b_0 + b_i$,
and its publication waits for that and for $\expub{0}$ at $b_0 + p_0$, giving
$b_0 + \max(b_i, p_0) + p_i$. Under \code{none}, its build ends at $b_i$ and its
publication waits for that and for $b_0 + p_0$. Fan-in. Under \code{publish},
$c$'s build starts when the last provider has published, at $\max_i (b_i +
p_i)$. Under \code{build}, it starts at $\max_i b_i$ and ends $b_c$ later, and
$c$'s publication waits for that and for every provider's publication. Under
\code{none}, $c$'s build ends at $b_c$ and only its publication waits for the
providers. Each is the longest path of \cref{prop:exec-longest}.
\end{proof}

\begin{proposition}[one publication lane, \spec{13.11}]
\label{prop:exec-lane}
Where one owner's publications are serialised (\spec{27.2}, \spec{28.6}) in
the order $1, \dots, n$ of \cref{def:order}, let $e_k$ be the instant the build
of package $k$ ends. Under every relation the run lasts
\[
  \max_{1 \le k \le n}\Bigl(e_k + \sum_{i=k}^{n} p_i\Bigr),
\]
which is at least $\sum_{i} p_i$ whatever the relation, and a weaker relation
only lowers the $e_k$.
\end{proposition}

\begin{proof}
Publication $k$ starts at $\max(e_k, E_{k-1})$, since it waits for its own
build and for the lane, and the lane order agrees with \cref{def:order}, so
every provider of $k$ has published by $E_{k-1}$. The recurrence is that of
\cref{prop:exec-chain} with $e_k$ in place of $B_k$, and the same induction
gives the maximum. The term $k = 1$ is $e_1 + \sum_i p_i \ge \sum_i p_i$. For
the last claim, add the lane edges $\expub{k-1} \to \expub{k}$ to the task
graph; they follow the publish order, which is a topological order of the
workspace graph, so the key of \cref{prop:readiness-acyclic} still increases
along every edge and the graph stays acyclic. With a fixed lane order every
task starts at the greatest finish of its predecessors, and by induction along
a topological order each finish time under a subset of the edges is at most
the finish time under the full set; $e_k$ is such a finish time.
\end{proof}

\begin{proposition}[budgets, \spec{13.11}]
\label{prop:exec-budget}
Let there be $m_b$ build slots and $m_p$ publication slots, and write $W_b =
\sum_i b_i$, $W_p = \sum_i p_i$ and $L$ for the longest path of the task
graph. No schedule finishes before $\max(L,\ W_b/m_b,\ W_p/m_p)$. A schedule
that never leaves a slot of a stage idle while a task of that stage is ready
finishes within $L + W_b/m_b + W_p/m_p$, which is within three times the best
possible. Under budgets a weaker relation can lengthen such a schedule, but
never the best one.
\end{proposition}

\begin{proof}
The lower bound: a schedule contains every path, so it lasts at least $L$; it
performs $W_b$ units of build work on $m_b$ slots, so at least $W_b/m_b$; and
likewise $W_p/m_p$. The upper bound: let the schedule end at $T$ and let $t_1$
be a task finishing at $T$. For $j \ge 1$, while $t_j$ has predecessors, let
$t_{j+1}$ be one that finishes last among them; this gives a chain $t_q \to
\dots \to t_1$ whose total duration is at most $L$. Consider an instant of
$[0, T)$ at which no task of the chain is running. It lies between the finish
of some $t_{j+1}$ and the start of $t_j$ (or before the start of $t_q$), and
at that instant $t_j$ is ready, since its latest predecessor has finished, and
has not started; by hypothesis every slot of $t_j$'s stage is then busy. The
instants at which all $m_b$ build slots are busy total at most $W_b/m_b$, and
those at which all $m_p$ publication slots are busy at most $W_p/m_p$. Hence $T
\le L + W_b/m_b + W_p/m_p$, and each summand is at most the lower bound. For
the last claim, a schedule feasible under a stronger relation is feasible under
a weaker one, so the best weaker schedule is no longer; that a particular
greedy schedule can be longer is \cref{ex:exec-anomaly}.
\end{proof}

The bound is Graham's argument for list scheduling under precedence
constraints~\cite{graham66,graham69}, applied to two stages with separate
budgets; \spec{13.11} does not claim it tight for this shape of graph.

\begin{example}[Graham's anomaly, \spec{13.11}]
\label{ex:exec-anomaly}
Take two build slots and one publication slot, packages $A = (2, 3)$ and $B =
(1, 2)$ where $B$ consumes $A$, and an unrelated $C = (8, 3)$, with ready tasks
taken in the order $A, B, C$. Under \code{build}, $B$'s build is not ready at
$0$: $A$ builds in $[0,2)$ and $C$ in $[0,8)$ on the two slots; $B$ builds in
$[2,3)$; $A$ publishes in $[2,5)$, $B$ in $[5,7)$ once $A$ has published, and
$C$ in $[8,11)$ once its build ends. The run lasts $11$. Under \code{none},
$B$'s build is ready at $0$ and takes the second slot: $A$ builds in $[0,2)$
and $B$ in $[0,1)$, so $C$ builds only in $[1,9)$; $A$ publishes in $[2,5)$,
$B$ in $[5,7)$ and $C$ in $[9,12)$. The run lasts $12$: the consumer's early
build took the slot the long build would have had. A list that preferred $C$
to $B$ at $0$ would finish at $11$ under \code{none} as well, which is the
point: the length depends on the list, not on the relation, and it is why
\spec{19.2a} promises a happens-before relation and no duration.
\end{example}

\begin{example}[an infrastructure chain]
\label{ex:exec-chain}
Let the infrastructure leg of \cref{ex:terraform} be consumed by a backend
consumed in turn by a frontend, with $(b, p)$ in minutes of $(4, 6)$ for the
infrastructure, $(9, 3)$ for the backend and $(12, 2)$ for the frontend. By
\cref{prop:exec-chain}, under \code{publish} the run lasts $10 + 12 + 14 = 36$
minutes; under \code{build} it lasts $\max(4 + 11,\ 13 + 5,\ 25 + 2) = 27$;
under \code{none} it lasts $\max(4 + 11,\ 9 + 5,\ 12 + 2) = 15$. The
publications still run in order in every case, and the frontend's publication
still ends the run.
\end{example}
